\section{Related work}\label{sec:related}

\paragraph{Interchange interventions and causal abstraction.}
Activation patching and interchange interventions test whether internal states carry the information an
explanation assigns to them \citep{vig2020gender,meng2022rome,wang2023ioi,geiger2021causal}. Causal abstraction
makes the claim precise: a high-level model abstracts a network if the two agree under interventions on aligned
variables, including interventions on several variables at once \citep{beckers2019abstracting,
rubenstein2017consistency,geiger2025abstraction}. Whether an abstraction is exact depends on the set of
interventions considered \citep{rubenstein2017consistency}, and a test of one variable at a time certifies a
claim about that variable with the others marginalised \citep{geiger2022iit}. DAS learns the alignment
\citep{geiger2024das,wu2023boundless}; its original training mixes single and two-variable interchanges
\citep[Appendix~A.1]{geiger2024das}, and alignments are usually scored by single-variable interchange accuracy, as in
the causal-variable track of MIB \citep{mueller2025mib} and the cause and isolation scores of RAVEL
\citep{huang2024ravel}. RAVEL scores an alignment both by whether an intervention changes the target attribute and by
whether it leaves the other attributes alone; the single queries in which the swapped variable keeps its value, which
turn out to matter in Study~4, are a test of the second kind. MIB's IOI model assigns the position variable to heads
7.3, 7.9 and 8.6 and the token variable to head 8.10 of GPT-2 Small; Study~4 tests that assignment under our
counterfactual model rather than under MIB's own scoring. For orthogonal linear alignments at one site,
\cref{prop:redundant} says when the multi-variable requirement adds nothing beyond single-variable tests over the full
domain. \citet{geiger2021causal} remark informally that single interchanges establish the abstraction of their
example, whose two variables occupy different neurons of one hidden layer and are therefore orthogonal subspaces of
one site; \cref{prop:redundant} states the condition under which this holds.

\paragraph{Explanations beyond the inputs used to find them.}
An explanation that holds on the inputs used to find or check it can fail elsewhere. \citet{bolukbasi2021illusion}
showed this for neuron interpretations across corpora, and \citet{friedman2024illusions} for simplified models
fitted on one distribution and tested on systematic-generalisation splits, the same mechanism as the restricted
alignments of Study~1. \citet{rai2026datadriven} find circuits with high in-distribution and low cross-dataset
faithfulness, and Experiment~5 of \citet{meloux2026validating} contains an invalid abstraction that differs from
the valid one only at a mediator value no natural input reaches. Learned alignments can also fail on held-out
entities \citep{huang2024ravel}, and an alignment map flexible enough to fit a randomly initialised network fails on
held-out names \citep{sutter2025nonlinear}. Failures of an abstraction can also be confined to part of the input
space: \citet{puyin2026bucketing} split the inputs by their pairwise interchange behaviour into regions where an
interpretation holds and regions where it fails, and \citet{pislar2025combining} combine simple causal models, each
explaining part of the inputs, trading the share of inputs explained against interchange accuracy.
\citet{heimersheim2024patching} recommend patching experiments whose inputs vary every hypothesised variable, and
\citet{wu2024reply} testing on examples outside the training data. Our coverage results replicate this effect under
exact labels, the held-out split of Study~2 asks how much it matters for single and joint tests, and Study~4's
confirmation uses names and templates that development never saw.

\paragraph{Interventions that leave the natural states.}
\citet{zhang2024patching} note that corrupting inputs with noise moves internal states off their natural
distribution. \citet{makelov2024subspace} show that a subspace patch can change the output through a dormant
pathway, and \citet{wu2024reply} that faithful subspaces generally overlap directions the downstream computation
ignores. \citet{grant2026divergent} distinguish harmless divergence from pernicious divergence, which uses hidden
pathways or causes dormant behavioural changes that appear only in other contexts, and propose checking interventions
in other contexts; our transported-divergence panel follows their Appendix~A.3, and their characterisation of sets
closed under coordinate patching (Theorem~A.2) is the closest formal precedent for \cref{prop:redundant}.
\Cref{cor:irreducible} connects joint tests to this work: they can only expose single edits of both variables whose
effect is dormant in a single test.

\paragraph{Interactions between interventions.}
The closest theory to \cref{prop:margin} is that of \citet{vaidyanathan2026mediators}. They show that the effect
activation patching attributes to a component is its pure effect plus an interaction term, which measures how far
that effect depends on the state of the rest of the network. To second order, the interaction term is the mixed
second derivative of the downstream map along the two displacements that the change of input causes, that of the
component and that of the rest of the network, so this term vanishes where the map is locally affine. Their setting
is the patching of whole components, such as attention heads at different sites, rather than of subspaces aligned
with the variables of a high-level causal model. The conditions also differ: their interaction term needs curvature
as well as displacement, whereas \cref{prop:margin} compares a joint query with a single query it reduces to, so that
a first-order drop of margin is enough, provided the other edit's miss is large compared with that query's margin.
\Cref{sec:compiled} discusses their remark that compiled networks are piecewise affine.
\citet{erramilli2026tomography} treats interventions as designed measurements of a model's effects and interactions;
in GPT-2 Small on the indirect-object task, the interaction between name-mover and negative name-mover heads is the
largest held-out term of the three pairs of head groups tested.

\paragraph{Evaluating validation tests against known answers.}
Compiled and trained networks with known circuits are used to evaluate interpretability methods
\citep{lindner2023tracr,conmy2023acdc,gupta2024interpbench,thurnherr2024tracrbench}. \citet{meloux2026validating}
benchmark metrics for causal abstraction on ten simulated systems with known valid and invalid explanations,
including a Tracr system; their metrics sample interventions of several cardinalities, the faithfulness-aware ones
also overwrite unmapped variables with sampled noise, and they argue that interventions on several nodes catch
interactions that single checks miss. The complement resampling of Study~1 follows causal scrubbing
\citep{chan2022scrubbing}. Several distinct explanations can fit the same network
\citep{meloux2025identifiable,chen2026rome}, which is why our labels concern one precise claim rather than the
uniqueness of any circuit. We reran two published demonstrations as checks (\cref{app:reproductions}).
